\documentclass[10pt]{article} 
\renewcommand{\thempfootnote}{\arabic{mpfootnote}}

\usepackage[hidelinks]{hyperref}

\usepackage[utf8]{inputenc}

\usepackage[T1]{fontenc}
\usepackage[ttdefault=true]{AnonymousPro}

\usepackage{multicol}
\setlength{\columnsep}{1cm}

\usepackage{abc}
\usepackage{graphicx}
\usepackage{musicography}

\usepackage[margin=0.5in]{geometry} 
\geometry{letterpaper}

\usepackage{fancyhdr} 
\pagestyle{fancy} % options: empty , plain , fancy
\renewcommand{\headrulewidth}{0pt} % customise the layout...
\lhead{}\chead{}\rhead{}
\lfoot{}\cfoot{\thepage}\rfoot{}

\begin{document}

\thispagestyle{empty}

\begin{abc}[name=detroitschottische,options={-O= +c -F single.fmt}]
X:1
T:Detroit Schottische
C:transcribed from Vivian Williams' version on Dance Music on the Oregon Trail album
C:by Kitty Steetle
M:4/4
L:1/8
K:D
f2 f>e d2 d>B | A>B A>F D2 F>D | E2 E>F G2 F>E | D>E F>G A4 | 
f2 f>e d2 d>B | A>B A>F D2 F>D | E>D E>F G>A B>c | d2 d2 !fine!d4 :|
K:A
e2 c2 e2 c2 | e>f e>c A2 c>A | B2 B>c d2 c>B | A>B c>d e4 | 
e>f e>c e>f e>c | e>f e>c A2 c>A | B>A B>c d>e f>g |1 a2 a2 a4 :|2  a4 !D.C.alfine! =g4|
\end{abc}

\vskip 0.5in

\noindent\makebox[\textwidth][c]{
  \begin{minipage}{7.0in}
    \setlength{\parindent}{1.5em}


    \textbf{Detroit Schottische} is attributed to Adam Couse, a
    dancing master who also owned a music store in Detroit, Michigan
    in the mid-19th
    century.\footnote{\url{https://tunearch.org/wiki/Annotation:Detroit_Schottische}}\footnote{Notably,
      Couse became friends with composer Stephen Foster who dedicated
      the ballad, \textbf{The Voice of By Gone Days}, to Couse.} Adam
    Couse’s Detroit Schottisch\footnote{No 'e'.} was published in 1854 as a piano tune
    with accompaniment and the facsimile can be found on the Library
    of Congress website.\footnote{\url{https://www.loc.gov/resource/music.mussm2-sm1854-230240/}}
    
    Elias Howe includes the tune in his Musician’s Omnibus No 1, under
    the section, May Flower Schottische Quadrilles, No 3 Detroit
    Schottische (1863, p 56).\footnote{You can download the entire
      Omnibus at the IMSLP link below. (Snip is taken from the online
      facsimile.)} Here Howe does not give a composer; although this
    is not wholly irregular for the time period, Howe does use the
    phrase ‘by permission’ for other tunes signifying that they have
    known composers (see Musician’s Omnibus No 1, \textbf{Dixie’s
      Land} and \textbf{Mocking Bird} both on page 8) leading one to
    think Howe believed the tune was
    traditional.\footnote{\url{https://imslp.org/wiki/Musician's_Omnibus_(Howe\%2C_Elias)}}
    
    The Detroit Schottische is very similar to \textbf{Flop-Eared Mule} and \textbf{Lop-Eared Mule} with both latter tunes being more popular in the American southeast. There is a strong likelihood that all three share a common ancestor and have no clear single source.\footnote{Sorry, Adam Couse.}
    
---Kitty Steetle


  \end{minipage}
}

\end{document}
